\section{Attack D: Shai-Hulud 2.0}
\label{sec:attackD}

Only two months after the Shai-Hulud 1.0 attack \cref{sec:attackC}, 
on November 24, 2025, a second, more aggressive attack compared to the previous one occurred, 
renamed Shai-Hulud 2.0, which this time infected 
1,055 \ac{NPM} packages within hours \cite{cutolo2025shaihulud20}.

The objective remained unchanged, namely 
to steal sensitive information such as tokens and authentication keys,
and the self-propagation also remained the same.

However, this new attack shows differences compared to its previous version \cref{sec:attackC} \cite{helixguard2025malicioussha1hulud}.

The new versions of these packages, published in the NPM registry, 
claimed to introduce the Node.js alternative Bun runtime,
inserting into the tarball a new installation file setup\_\allowbreak bun.js
and a file bun\_\allowbreak environment.js.

The first file mimics the code related to the Bun configuration and 
calls the second file bun\_\allowbreak environment.js.

The second file, weighing over 10 MB, 
contains a single obfuscated line of code 10,157,373 characters long.

As in the previous version \cref{sec:attackC}, 
once the script (bun\_\allowbreak environment.js) is started,
it downloads and executes a version of TruffleHog
with the purpose of stealing sensitive information such as \ac{NPM} tokens, 
\ac{AWS}, \ac{GCP}, and Azure credentials, and environment variables.

After finishing the search, the malicious code publishes the 
collected secrets and system information encoded
by creating a public repository 
on behalf of the victim named Sha1-Hulud: The Second Coming,
confirming that it is the successor of the previous Shai-Hulud 1.0 attack \cref{sec:attackC}.

In this new version of the attack, however, the malicious code
makes use of an obfuscation technique that encodes 
in hexadecimal function names, variables, and constants,
making the bun\_\allowbreak environment.js file full of hexadecimal values.

Even though the code is obfuscated, 
words like "password" and "aws\_secret" are visible 
and present in clear text, making it possible to identify them through static analysis.